




\section{FLOPs Comparison}
\label{sec:flops}




\begin{table*}[t]
\centering
\setlength{\tabcolsep}{6pt}
\newcommand{\rel}[1]{{\scriptsize\,(#1$\times$)}}
\resizebox{0.8\linewidth}{!}{
\begin{tabular}{@{}lcccc@{\hskip 14pt}ccc@{}}
\toprule
& \multicolumn{3}{c}{Inference (FLOPs/token, $\times10^{9}$)}
& \multicolumn{3}{c}{Training (total FLOPs, $\times10^{19}$)} \\
\cmidrule(r){2-4}\cmidrule(l){5-7}
$T$ & \textsc{LoopMTP} & LoopFormer & $\Delta$
    & \textsc{LoopMTP} & LoopFormer & $\Delta$ \\
\midrule
Non-looped & \multicolumn{2}{c}{0.481} & --
           & \multicolumn{2}{c}{0.988} & -- \\
\midrule
3 & 1.20\rel{2.49} & 1.31\rel{2.73} & $-8.8\%$ & 2.46\rel{2.49} & 4.15\rel{4.20} & $-40.7\%$ \\
5 & 1.93\rel{4.01} & 2.12\rel{4.40} & $-8.9\%$ & 3.96\rel{4.01} & 6.63\rel{6.71} & $-40.2\%$ \\
7 & 2.66\rel{5.53} & 2.92\rel{6.08} & $-9.0\%$ & 5.47\rel{5.53} & 9.11\rel{9.23} & $-40.1\%$ \\
9 & 3.39\rel{7.05} & 3.73\rel{7.75} & $-9.0\%$ & 6.97\rel{7.05} & 11.6\rel{11.74} & $-39.9\%$ \\
\bottomrule
\end{tabular}
}
\caption{\textbf{FLOPs comparison between \textsc{LoopMTP}, its non-looped baseline and LoopFormer \cite{jeddi2026loopformer}.} Parenthesised values are multiples of the
non-looped baseline. $\Delta$ is the relative change of \textsc{LoopMTP}
with respect to LoopFormer; negative values indicate that
\textsc{LoopMTP} requires less computation.}
\label{tab:flops-loopmtp-vs-loopformer}
\end{table*}



Following the accounting of the LoopFormer \cite{jeddi2026loopformer}, we write the cost of a
forward pass with $t$ loop iterations as $C(t) = C_{io} + t\, C_1$, where
$C_{io}$ collects the loop-independent input/output compute (final norm + unembedding;
embedding lookups are free) and $C_1$ is the cost of one pass through the shared
$L$-block stack plus each model's per-iteration machinery. 
The non-looped baseline is the $T=1$ instance of the same expression, with a slightly wider \textsc{SwiGLU}. \textsc{LoopMTP} follows the LLaMA sizing rule ($\tfrac{2}{3}\cdot 4d$, rounded up to the nearest multiple of $256$), so the nominal $4096$ of Table~\ref{tab:hyperparams_config} corresponds to a realized hidden width of $2816$ at $d{=}1024$. For the baseline, which has no input projection or aggregation gate, we widen this to the next multiple of $256$, i.e.\ $3072$. Because the width is constrained to multiples of $256$, this does not match the recurrence parameters exactly: it overshoots them, leaving the baseline roughly $6$M parameters \emph{larger} than \textsc{LoopMTP}.
We count $2$ FLOPs per multiply--accumulate, score attention
causally, and take one training step to cost $3\times$ the forward pass.
For our configuration ($L=12$, $d=1024$, $S=2048$, $|V|=50304$, $D \approx 6.8$B), per token:
\begin{align*}
\text{Baseline:}\quad &
\left\{
\begin{aligned}
    C_{io}/S &= 1.03 \times 10^{8}\\ 
    C_1^{\mathrm{base}}/S &= 3.78 \times 10^{8}
\end{aligned}
\right. \\[7pt]
\text{LoopMTP:}\quad &
\left\{
\begin{aligned}
    C_{io}/S &= 1.03 \times 10^{8}\\
    C_1/S &= 3.65 \times 10^{8}
\end{aligned}
\right. \\[7pt]
\text{LoopFormer:}\quad &
\left\{
\begin{aligned}
    C_{io}/S &= 1.03 \times 10^{8}\\
    C_1/S &= 4.03 \times 10^{8}
\end{aligned}
\right.
\end{align*}

Inference uses a single $T$-loop trajectory for both models,
$C_{\mathrm{inf}}(T) = C_{io} + T\,C_1$ and a single pass for the baseline,
$C_{\mathrm{inf}}^{\mathrm{base}} = C_{io} + C_1^{\mathrm{base}}$ . Over a token budget $D$, training costs:
\begin{align*}
\mathcal{C}^{\mathrm{train}}_{\mathrm{base}} &= 3\,\big(C_{io} + C_1^{\mathrm{base}}\big)\, D/S, \\[8pt]
\mathcal{C}^{\mathrm{train}}_{\mathrm{LoopMTP}} &= 3\,\big(C_{io} + T\, C_1\big)\, D/S, \\[8pt]
\mathbb{E}\big[\mathcal{C}^{\mathrm{train}}_{\mathrm{LoopFormer}}\big] &= 3\,\big(2\,C_{io} + \tfrac{3T}{2}\, C_1\big)\, D/S,
\end{align*}
since LoopFormer backpropagates both its full $T$-loop trajectory and a short
$M$-loop trajectory with $M \sim \mathrm{Unif}\{1,\dots,T-1\}$ ($\mathbb{E}[M] = T/2$).


Table~\ref{tab:flops-loopmtp-vs-loopformer} compares the training and inference FLOPs of LoopFormer and \textsc{LoopMTP}. For a matched loop count $T$ and token budget, \textsc{LoopMTP} requires only 0.59--0.60$\times$ the training FLOPs of LoopFormer, while achieving better performance. The additional computational cost of LoopFormer primarily stems from its dual-trajectory objective. During inference, \textsc{LoopMTP} remains approximately 10\% more FLOP-efficient.
Relative to the non-looped baseline of the same size, \textsc{LoopMTP} costs 2.49--7.05$\times$ the FLOPs at both training and inference for $T{=}3$--$9$. Importantly, looping converts compute into effective depth; it does not reduce total computation relative to a shallow model.

